\subsection{Compound Inequalities} 
Sometimes, we need a quantity to fit \emph{more than one} condition at the same time.
\begin{Example}
A class 2a light truck weighs over 6000 points, but not more than 8500 pounds. If the variable $w$ represents the weight of a truck, we would describe these restrictions with
\[
w>6000
\ \makeblanksmall\ 
w\leq 8500
\]
\end{Example}
Sometimes, we need a quantity to fit \emph{at least one} of several conditions.
\begin{Example}
An online printing company offers short-run prints of 100 or fewer units, or large runs of 2000 or more units. If the variable $n$ represents the number of copies a user can order, we would describe these restrictions with
\[
n\leq 100
\ \makeblanksmall\ 
n\geq 2000
\]
\end{Example}
\begin{definition}[Compound Inequality]
A \term{compound inequality} is a collection of two or more inequalities, combined using the word ``and'' or the word ``or''.
\begin{itemize}
\item If they are combined with ``and,'' the solution set is the \emph{intersection} of the two inequalities' solution sets.
\item If they are combined with ``or,'' the solution set is the \emph{union} of the two inequalities' solution sets.
\end{itemize}
An inequality that is not compound is a \term{simple inequality}.
\end{definition}
\begin{Example}
Sketch the solution set of
\[
x\leq 5
\mbox{ and }
x>-2
\]
on the number line below.
\begin{icpic}
\useasboundingbox (-6,-1) rectangle (8,1);
\draw[axis] (-6,0) -- (8,0);
\foreach \x in {0,...,7} \draw[tick] (\x,0.2) -- (\x,-0.2) node[anchor=north]{$\x$};
\foreach \x in {1,...,5} \draw[tick] (-\x,0.2) -- (-\x,-0.2) node[anchor=north]{$-\x\ $};
\end{icpic}
\end{Example}
\begin{Example}
Sketch the solution set of
\[
x<-4
\mbox{ or }
x\geq 3
\]
on the number line below.
\begin{icpic}
\useasboundingbox (-6,-1) rectangle (8,1);
\draw[axis] (-7,0) -- (7,0);
\foreach \x in {0,...,6} \draw[tick] (\x,0.2) -- (\x,-0.2) node[anchor=north]{$\x$};
\foreach \x in {1,...,6} \draw[tick] (-\x,0.2) -- (-\x,-0.2) node[anchor=north]{$-\x\ $};
\end{icpic}
\end{Example}